\section{Discussion}
\label{sec:discussion}

\subsection{Implications of Planetary Magnetic Environments for Human Habitats}
The central insight established by our integrated analysis is that space biological systems and human exploration habitats deployed to the Moon and Mars will operate within two fundamentally distinct magnetic operational paradigms:

\begin{enumerate}
    \item \textbf{The Homogeneous Near-Null Lunar \& Martian Lowland Paradigm ($<\SI{25}{\nano\tesla}$):} Primary candidate sites for initial crewed bases—specifically all 13 Artemis III South Polar regions on the Moon ($< \SI{1.0}{\nano\tesla}$) and ice-rich ISRU candidate sites on Mars such as Arcadia Planitia and Deuteronilus Mensae ($< \SI{20}{\nano\tesla}$ surface)—are selected based on volatile resource availability (water ice) and solar illumination. However, our findings demonstrate that these sites represent extreme near-null hypomagnetic environments ($>2,500\times$ to $>60,000\times$ weaker than Earth's GMF). Bioregenerative life support systems (BLSS), agricultural growth chambers, and human crew residing in these habitats will undergo chronic, unmitigated magnetic deprivation.
    \item \textbf{The Martian Southern Highland Mini-Magnetosphere Outposts ($>\SI{1500}{\nano\tesla}$):} In contrast, specialized scientific outposts located within the Terra Sirenum or Terra Cimmeria anomaly belts on Mars reside within intense remnant magnetic fields ($>\SI{1500}{\nano\tesla}$ surface). These fields form mini-magnetospheres extending hundreds of kilometers into the Martian ionosphere \citep{Mitchell2001}, capable of partially deflecting low-energy solar wind ions and creating localized plasma cusps. While these fields remain $\sim 25\times$ to $50\times$ below Earth's baseline, they provide a substantially elevated baseline compared to northern lowland sites.
\end{enumerate}

\subsection{The Triple-Stressor Triad: Hypomagnetism, Altered Gravity, and Radiation}
On extraterrestrial surfaces, magnetic field deprivation will not act in isolation. Instead, biological systems will confront a unique triple-stressor triad:
\begin{itemize}
    \item Chronic hypomagnetic field deprivation ($< \SI{1.0}{\nano\tesla}$ on Moon; $< \SI{20}{\nano\tesla}$ on Martian northern plains)
    \item Reduced partial gravity ($0.166\,g$ on Moon; $0.38\,g$ on Mars)
    \item Chronic galactic cosmic radiation (GCR) and solar particle events (SPE), delivering surface doses of $\sim 60\text{ to }\SI{100}{\milli\sievert}/\text{year}$ on Moon and $\sim 200\text{ to }\SI{300}{\milli\sievert}/\text{year}$ on Mars.
\end{itemize}

Crucially, experimental evidence indicates potent biochemical synergies among these environmental factors. In mammalian hindlimb unloading models, HMF exposure significantly aggravates musculoskeletal bone loss and trabecular demineralization \citep{Xu2012, Mo2014}. Furthermore, HMF-induced perturbations in radical pair singlet-triplet spin-state recombination kinetics and ROS accumulation \citep{Binhi2017, Hore2016} can impair cellular DNA repair pathways, potentially sensitizing biological systems to ionizing radiation damage.

\subsection{Biophysical Mechanisms of Biological Hypomagnetic Stress}
The biological sensitivity to hypomagnetic fields is mediated by two primary biophysical mechanisms:
\begin{enumerate}
    \item \textbf{The Radical Pair Mechanism (RPM):} Flavin-based photoreceptors (cryptochromes CRY1/CRY2) rely on light-induced radical pair recombination dynamics $[\text{FAD}^{\bullet-}, \text{TrpH}^{\bullet+}]$. Removal of the background geomagnetic field alters the singlet-triplet spin interconversion rate, perturbing downstream signaling pathways governing floral transition in plants \citep{Agliassa2018} and circadian clock entrainment in animals \citep{Yoshii2009}.
    \item \textbf{Cellular Redox and Auxin Efflux Transport:} Shielding the geomagnetic field disorganizes thylakoid membrane ultrastructure \citep{Belyavskaya2004} and elevates mitochondrial ROS generation \citep{Maffei2014}. In plant roots, this redox imbalance disrupts the polar localization of PIN2 efflux carriers, impairing directional auxin distribution and gravitropic bending \citep{Narayana2018}.
\end{enumerate}

\subsection{Operational Recommendations for Lunar and Martian Surface Operations}
To de-risk future planetary exploration and ensure the viability of long-duration surface life support systems, we propose four core operational recommendations:
\begin{enumerate}
    \item \textbf{Standardized Surface Magnetometry Integration:} All future precursor landers, rovers, and human surface payloads (Artemis III, CLPS, Commercial Mars missions) should integrate calibrated fluxgate magnetometers to record high-resolution ground-level magnetic fields at human operational scale.
    \item \textbf{Active Magnetic Field Restoration in Growth Modules:} Agricultural and BLSS modules deployed in lunar polar craters or Martian northern lowlands should incorporate active Helmholtz coil arrays to generate artificial Earth-like static magnetic fields ($\sim\SI{50}{\micro\tesla}$), restoring normative cryptochrome signaling, auxin transport, and crop yields.
    \item \textbf{Controlled Magnetic-Biology Payloads:} Early surface missions should deploy dual-chamber biological experiments (shielded near-null vs. \SI{50}{\micro\tesla} artificial field) directly on the lunar and Martian surface to empirically isolate hypomagnetic effects from radiation and partial gravity.
    \item \textbf{Habitat Siting Optimization:} Long-term infrastructure planning should incorporate magnetic field architecture alongside volatile access, illumination, and slope topography when evaluating permanent base siting.
\end{enumerate}
